\chapter{Buổi 3}
\label{Chapter3}
\setcounter{secnumdepth}{5}

\section*{Câu 1}
\addcontentsline{toc}{section}{Câu 1}
Ứng dụng xử lý ảnh và video trong thương mại điện tử và tiếp thị số giúp ích gì?

Xử lý ảnh và video giúp doanh nghiệp hiểu sản phẩm, hành vi khách hàng và ngữ cảnh mua sắm tốt hơn. Nhờ đó, hệ thống có thể quảng bá, tìm kiếm và gợi ý sản phẩm chính xác hơn.

\begin{itemize}
    \item \textbf{Tăng khả năng nhận diện sản phẩm:} Hệ thống có thể phát hiện logo, bao bì, kiểu dáng hoặc màu sắc sản phẩm trong ảnh và video quảng cáo.

    \item \textbf{Hỗ trợ tìm kiếm sản phẩm:} Người dùng có thể tải ảnh lên để tìm sản phẩm tương tự, thay vì phải mô tả bằng từ khóa.

    \item \textbf{Cá nhân hóa gợi ý:} Mô hình có thể phân tích phong cách, màu sắc và loại sản phẩm khách hàng quan tâm để đề xuất sản phẩm phù hợp hơn.

    \item \textbf{Tối ưu tiếp thị số:} Phân tích hình ảnh, video và phản ứng của người dùng giúp doanh nghiệp chọn nội dung quảng cáo hấp dẫn, đúng nhóm khách hàng mục tiêu.
\end{itemize}


\section*{Câu 2}
\addcontentsline{toc}{section}{Câu 2}
Bài toán nhận dạng mẫu và truy vấn mẫu hỗ trợ đưa thông tin sản phẩm đến khách hàng như thế nào?

Nhận dạng mẫu giúp hệ thống tự động phát hiện sản phẩm, logo, kiểu dáng hoặc đặc điểm thị giác trong ảnh và video. Từ đó, hệ thống có thể gắn nhãn sản phẩm, phân loại danh mục và hiển thị thông tin phù hợp cho khách hàng.

Truy vấn mẫu cho phép khách hàng dùng ảnh, video hoặc mô tả để tìm sản phẩm tương tự. Ví dụ, người dùng chụp ảnh một đôi giày hoặc chiếc áo, hệ thống sẽ so sánh đặc trưng hình ảnh và đề xuất các sản phẩm gần giống trong kho hàng.

Nhờ hai bài toán này, thông tin sản phẩm được truyền đến khách hàng nhanh hơn, đúng nhu cầu hơn và ít phụ thuộc vào việc người dùng phải nhập từ khóa chính xác.


\section*{Câu 3}
\addcontentsline{toc}{section}{Câu 3}
Ý nghĩa của các ứng dụng giám sát, quản lý và đếm người trong quản lý xã hội là gì?

Các hệ thống giám sát, quản lý và đếm người giúp cơ quan quản lý hiểu tình trạng hoạt động của không gian công cộng theo thời gian thực. Thay vì chỉ dựa vào quan sát thủ công, dữ liệu từ camera có thể hỗ trợ ra quyết định nhanh và chính xác hơn.

\begin{itemize}
    \item \textbf{Quản lý an ninh:} Phát hiện đám đông bất thường, hành vi nguy hiểm, xâm nhập trái phép hoặc tình huống cần can thiệp khẩn cấp.

    \item \textbf{Điều phối hạ tầng:} Đếm người tại nhà ga, sân bay, bệnh viện hoặc trung tâm thương mại để điều chỉnh luồng di chuyển và giảm ùn tắc.

    \item \textbf{Hỗ trợ quy hoạch:} Dữ liệu mật độ người giúp đánh giá mức sử dụng không gian công cộng, từ đó cải thiện giao thông, dịch vụ và bố trí nhân lực.

    \item \textbf{Ứng phó khẩn cấp:} Khi xảy ra cháy nổ, thiên tai hoặc sự kiện đông người, hệ thống có thể ước lượng số người trong khu vực để hỗ trợ sơ tán và cứu hộ.
\end{itemize}


\section*{Câu 4}
\addcontentsline{toc}{section}{Câu 4}
Vai trò của Khoa học dữ liệu thị giác trong xây dựng thành phố thông minh.

Khoa học dữ liệu thị giác đóng vai trò quan trọng trong thành phố thông minh vì camera, cảm biến hình ảnh và video giám sát tạo ra lượng dữ liệu rất lớn mỗi ngày. Nếu chỉ lưu trữ thủ công, dữ liệu này khó khai thác hiệu quả. Khi kết hợp xử lý ảnh, thị giác máy tính và phân tích dữ liệu, hệ thống có thể tự động hiểu bối cảnh đô thị, phát hiện sự kiện và hỗ trợ cơ quan quản lý ra quyết định nhanh hơn.

\begin{itemize}
    \item \textbf{Truy tìm đối tượng:} Hệ thống có thể tìm người, phương tiện hoặc đồ vật nghi vấn trên nhiều camera khác nhau dựa vào khuôn mặt, quần áo, biển số, màu sắc hoặc hình dạng. Điều này giúp rút ngắn thời gian tìm kiếm trong các tình huống an ninh.

    \item \textbf{Phát hiện sự kiện bất thường:} Các mô hình thị giác có thể nhận biết hành vi đánh nhau, tụ tập đông bất thường, đi ngược chiều, xe dừng sai vị trí, khói lửa hoặc cháy nổ. Những sự kiện này được đánh dấu để nhân viên quản lý kiểm tra ngay.

    \item \textbf{Cảnh báo sớm:} Khi phát hiện dấu hiệu nguy hiểm, hệ thống có thể gửi cảnh báo theo thời gian thực đến trung tâm điều hành, lực lượng an ninh, cứu hỏa hoặc y tế. Nhờ đó, thành phố có thể phản ứng nhanh trước tai nạn, ùn tắc hoặc sự cố công cộng.

    \item \textbf{Truy vấn thông minh:} Người dùng có thể đặt câu hỏi hoặc nhập mô tả như ``người mặc áo đỏ xuất hiện lúc 8 giờ'' hoặc ``xe biển số lảng vảng gần cổng''. Hệ thống sẽ tìm trong dữ liệu video mà không cần con người xem lại toàn bộ camera.

    \item \textbf{Lưu trữ và khai thác dữ liệu hiệu quả:} Dữ liệu giám sát có thể được gắn nhãn theo thời gian, vị trí, đối tượng và sự kiện. Việc này giúp truy xuất nhanh, giảm chi phí kiểm tra thủ công và hỗ trợ điều tra sau sự cố.

    \item \textbf{Hỗ trợ quy hoạch đô thị:} Phân tích mật độ người, lưu lượng xe và điểm thường xảy ra ùn tắc giúp chính quyền tối ưu giao thông, bố trí nhân lực, cải thiện hạ tầng và nâng cao chất lượng dịch vụ công.
\end{itemize}


\section*{Câu 5}
\addcontentsline{toc}{section}{Câu 5}
So sánh giám sát bằng mắt người và giám sát bằng text query.

\small
\begin{longtable}{|p{0.22\textwidth}|p{0.34\textwidth}|p{0.34\textwidth}|}
\hline
\textbf{Tiêu chí} & \textbf{Giám sát bằng mắt người} & \textbf{Giám sát bằng text query} \\
\hline
\textbf{Cách hoạt động} & Con người quan sát camera liên tục và tự phát hiện sự kiện. & Người dùng nhập mô tả, ví dụ ``người mặc áo đỏ'' hoặc ``xe dừng trước cổng'', hệ thống tự truy vấn video. \\
\hline
\textbf{Khối lượng công việc} & Cao, vì phải theo dõi nhiều màn hình trong thời gian dài. & Thấp hơn, vì hệ thống tự lọc dữ liệu theo yêu cầu truy vấn. \\
\hline
\textbf{Phạm vi tìm kiếm} & Bị giới hạn bởi số camera và khả năng tập trung của người quan sát. & Có thể tìm trên nhiều camera, nhiều thời điểm và dữ liệu đã lưu trữ. \\
\hline
\textbf{Tốc độ phản hồi} & Phụ thuộc vào kinh nghiệm và mức độ tập trung của nhân viên. & Nhanh hơn khi dữ liệu đã được lập chỉ mục và mô hình nhận dạng hoạt động tốt. \\
\hline
\textbf{Hạn chế} & Dễ mệt mỏi, bỏ sót sự kiện nhỏ hoặc bất thường hiếm gặp. & Phụ thuộc vào chất lượng mô hình, độ rõ của mô tả và chất lượng video. \\
\hline
\end{longtable}
\normalsize


\section*{Câu 6}
\addcontentsline{toc}{section}{Câu 6}
Phát biểu bài toán Person-Reidentification.

Person-Reidentification là bài toán nhận diện lại cùng một người khi người đó xuất hiện ở các camera khác nhau, hoặc ở các thời điểm khác nhau trong cùng một hệ thống giám sát. Mục tiêu không chỉ là phát hiện có người trong ảnh, mà còn xác định liệu người trong ảnh truy vấn có trùng với người đã xuất hiện trong cơ sở dữ liệu hay không.

Có thể phát biểu bài toán như sau: cho một ảnh hoặc đoạn video truy vấn chứa một người cần tìm, hệ thống phải so sánh người đó với tập ảnh/video từ nhiều camera và trả về danh sách các ảnh có khả năng thuộc cùng một người, được sắp xếp theo độ tương đồng.

\begin{itemize}
    \item \textbf{Đầu vào:} Ảnh hoặc video truy vấn của một người, cùng với tập dữ liệu từ nhiều camera giám sát.

    \item \textbf{Đầu ra:} Danh sách các lần xuất hiện có khả năng là cùng một người, kèm vị trí camera, thời gian và độ tin cậy.

    \item \textbf{Đặc trưng sử dụng:} Màu sắc và kiểu quần áo, dáng người, balo, phụ kiện, kiểu tóc, tư thế hoặc đặc trưng học sâu trích xuất từ ảnh.

    \item \textbf{Khó khăn:} Góc nhìn camera khác nhau, ánh sáng thay đổi, che khuất, ảnh mờ, người đổi tư thế hoặc nhiều người có trang phục giống nhau.
\end{itemize}

Ví dụ, trong một trung tâm thương mại, nếu cần tìm một người mặc áo xanh xuất hiện ở cổng vào lúc 9 giờ, hệ thống ReID có thể truy vết người đó qua nhiều camera để biết họ đã đi qua khu vực nào và vào thời điểm nào. Bài toán này rất hữu ích trong giám sát an ninh, tìm người thất lạc và phân tích luồng di chuyển trong không gian công cộng.


\section*{Câu 7}
\addcontentsline{toc}{section}{Câu 7}
Phân biệt truy vấn sản phẩm và truy vấn người.

Truy vấn sản phẩm và truy vấn người đều là bài toán tìm kiếm dựa trên đặc trưng thị giác, nhưng mục tiêu và yêu cầu độ chính xác khác nhau.

\small
\begin{longtable}{|p{0.24\textwidth}|p{0.33\textwidth}|p{0.33\textwidth}|}
\hline
\textbf{Tiêu chí} & \textbf{Truy vấn sản phẩm} & \textbf{Truy vấn người} \\
\hline
\textbf{Mục tiêu} & Tìm sản phẩm giống hoặc \textbf{tương tự} về kiểu dáng, màu sắc, chất liệu, thương hiệu. & Tìm \textbf{đúng} người đã xuất hiện trong camera hoặc cơ sở dữ liệu. \\
\hline
\textbf{Yêu cầu chính xác} & Có thể chấp nhận kết quả gần giống, vì mục tiêu thường là gợi ý lựa chọn. & Cần độ chính xác cao, vì nhầm lẫn có thể ảnh hưởng đến an ninh hoặc quyền riêng tư. \\
\hline
\textbf{Đặc trưng quan trọng} & Hình dạng, màu sắc, logo, mẫu mã, chất liệu, danh mục sản phẩm. & Khuôn mặt, dáng người, quần áo, phụ kiện, tư thế, thời gian và vị trí xuất hiện. \\
\hline
\textbf{Rủi ro khi sai} & Thường chỉ làm giảm chất lượng gợi ý hoặc trải nghiệm mua sắm. & Có thể gây nhận diện nhầm, ảnh hưởng đến cá nhân hoặc quyết định điều tra. \\
\hline
\end{longtable}
\normalsize


\section*{Câu 8}
\addcontentsline{toc}{section}{Câu 8}
Phân loại bất thường đã biết trước và bất thường chưa biết trước như thế nào?

\small
\begin{longtable}{|p{0.25\textwidth}|p{0.32\textwidth}|p{0.33\textwidth}|}
\hline
\textbf{Tiêu chí} & \textbf{Bất thường đã biết trước} & \textbf{Bất thường chưa biết trước} \\
\hline
\textbf{Khái niệm} & Là các tình huống đã được định nghĩa rõ và thường có dữ liệu mẫu hoặc nhãn huấn luyện. & Là các tình huống lạ, hiếm hoặc chưa từng xuất hiện trong dữ liệu huấn luyện. \\
\hline
\textbf{Cách phát hiện} & Huấn luyện mô hình nhận diện trực tiếp từng loại bất thường cụ thể. & Học trạng thái bình thường, sau đó phát hiện những mẫu lệch khỏi quy luật bình thường. \\
\hline
\textbf{Ví dụ} & Cháy, khói, đánh nhau, té ngã, vượt rào, đi ngược chiều, xe dừng sai vị trí. & Vật thể lạ trong khu vực cấm, hành vi gây rối kiểu mới, tình huống hiếm hoặc cách gian lận chưa từng gặp. \\
\hline
\textbf{Dữ liệu cần có} & Cần dữ liệu có nhãn cho từng loại bất thường. & Chủ yếu cần dữ liệu bình thường; không nhất thiết phải có nhãn cho mọi loại bất thường. \\
\hline
\textbf{Ưu điểm} & Dễ đánh giá, dễ giải thích và cho kết quả tốt với các tình huống đã học. & Có khả năng phát hiện sự kiện mới, chưa được định nghĩa trước. \\
\hline
\textbf{Hạn chế} & Dễ bỏ sót nếu sự kiện xảy ra khác với mẫu đã huấn luyện. & Dễ báo động giả vì không phải mọi điểm khác thường đều nguy hiểm. \\
\hline
\end{longtable}
\normalsize

Ví dụ, trong camera giao thông, hành vi \textbf{đi ngược chiều} là bất thường đã biết trước vì có thể mô tả rõ và thu thập dữ liệu huấn luyện. Ngược lại, một vật thể lạ rơi xuống đường hoặc một kiểu gây rối mới trong đám đông có thể là bất thường chưa biết trước vì hệ thống chưa từng thấy trong dữ liệu.


\section*{Câu 9}
\addcontentsline{toc}{section}{Câu 9}
Triết lý để giải quyết bất thường đã biết trước và bất thường chưa biết trước.

Triết lý chung là: với bất thường đã biết trước, ta dạy hệ thống nhận diện trực tiếp các mẫu nguy hiểm đã được định nghĩa; với bất thường chưa biết trước, ta dạy hệ thống hiểu thế nào là bình thường để phát hiện những điểm lệch khỏi quy luật đó. Hai nhóm này cần cách tiếp cận khác nhau vì mức độ hiểu biết ban đầu về sự kiện là khác nhau.

\begin{itemize}
    \item \textbf{Bất thường đã biết trước:} Nên tiếp cận như bài toán học có giám sát. Ta thu thập dữ liệu có nhãn, ví dụ cháy, khói, đánh nhau, té ngã hoặc đi ngược chiều, rồi huấn luyện mô hình phát hiện trực tiếp từng loại sự kiện. Triết lý ở đây là ``biết nguy cơ nào thì mô hình hóa nguy cơ đó''.

    \item \textbf{Bất thường chưa biết trước:} Không thể liệt kê hết mọi tình huống có thể xảy ra, nên cần học mô hình của trạng thái bình thường. Khi một hành vi, vật thể hoặc chuyển động lệch mạnh khỏi quy luật bình thường, hệ thống sẽ đánh dấu là đáng nghi để con người kiểm tra. Triết lý ở đây là ``không biết trước mọi nguy cơ, nhưng có thể nhận ra điều không bình thường''.

    \item \textbf{Cập nhật liên tục:} Những bất thường chưa biết trước sau khi được xác minh có thể được đưa vào tập dữ liệu huấn luyện. Khi đó, hệ thống dần chuyển một phần ``chưa biết'' thành ``đã biết'', giúp mô hình ngày càng hoàn thiện.
\end{itemize}


\section*{Câu 10}
\addcontentsline{toc}{section}{Câu 10}
Vai trò của thiết bị tự hành.

Thiết bị tự hành là các hệ thống có khả năng tự cảm nhận môi trường, phân tích dữ liệu và thực hiện hành động mà không cần con người điều khiển trực tiếp liên tục. Ví dụ gồm xe tự hành, robot giao hàng, drone giám sát, robot kho hàng và robot trong nhà máy.

\begin{itemize}
    \item \textbf{Tăng hiệu quả vận hành:} Thiết bị tự hành có thể thực hiện các công việc lặp lại như vận chuyển hàng, kiểm kho, tuần tra, giao hàng hoặc giám sát dây chuyền sản xuất với tốc độ ổn định.

    \item \textbf{Giảm rủi ro cho con người:} Trong môi trường nguy hiểm như khu vực cháy nổ, công trường, nhà máy hóa chất, hầm mỏ hoặc vùng thiên tai, thiết bị tự hành có thể thay con người quan sát và thực hiện nhiệm vụ.

    \item \textbf{Mở rộng khả năng giám sát:} Drone hoặc robot tuần tra có thể di chuyển đến nhiều vị trí khác nhau, thu thập ảnh/video và phát hiện bất thường ở những khu vực camera cố định khó bao phủ.

    \item \textbf{Ra quyết định theo thời gian thực:} Nhờ cảm biến, xử lý ảnh, định vị và mô hình AI, thiết bị tự hành có thể phát hiện vật cản, chọn đường đi, tránh va chạm và điều chỉnh hành vi theo môi trường.

    \item \textbf{Hỗ trợ sản xuất và logistics:} Trong nhà máy hoặc kho hàng, robot tự hành giúp vận chuyển linh kiện, phân loại hàng hóa, giao hàng nội bộ và phối hợp với dây chuyền tự động để giảm thời gian chờ.
\end{itemize}


\section*{Câu 11}
\addcontentsline{toc}{section}{Câu 11}
Các tác vụ sau giúp gì cho thiết bị tự hành?

\begin{enumerate}
    \item \textbf{Ổn định chất lượng video:} Camera trên robot, drone hoặc xe tự hành thường bị rung, mờ, thay đổi ánh sáng hoặc nhiễu do chuyển động. Ổn định video giúp hình ảnh rõ hơn, giảm sai lệch khi phát hiện vật thể, nhận diện làn đường hoặc theo dõi mục tiêu.

    \item \textbf{Định vị tự động:} Thiết bị tự hành cần biết vị trí của mình trong môi trường. Các phương pháp như visual odometry hoặc SLAM dùng hình ảnh từ camera để ước lượng vị trí, hướng di chuyển và quỹ đạo, đặc biệt hữu ích khi GPS yếu hoặc không có GPS.

    \item \textbf{Tái tạo môi trường 3D:} Từ ảnh, video, LiDAR hoặc cảm biến độ sâu, hệ thống có thể xây dựng bản đồ 3D của môi trường. Bản đồ này giúp thiết bị hiểu khoảng cách, bề mặt, vật cản, hành lang di chuyển và không gian có thể đi qua.

    \item \textbf{Phát hiện và tránh vật cản:} Thiết bị tự hành phải nhận biết người, xe, tường, cột, hàng hóa hoặc vật thể bất ngờ xuất hiện. Sau khi phát hiện vật cản, hệ thống cần điều chỉnh tốc độ, dừng lại hoặc chọn đường đi khác để tránh va chạm.
\end{enumerate}


\section*{Câu 12}
\addcontentsline{toc}{section}{Câu 12}
Ứng dụng của thiết bị tự hành vào quy trình sản xuất và quản lý xã hội.

\begin{itemize}
    \item \textbf{Nông nghiệp thông minh:} Drone hoặc robot tự hành có thể phun thuốc trừ sâu, tưới nước, gieo hạt, theo dõi sâu bệnh và đánh giá tình trạng cây trồng bằng camera RGB, hồng ngoại hoặc đa phổ. Nhờ đó, việc canh tác chính xác hơn và giảm lãng phí hóa chất.

    \item \textbf{Chăn nuôi:} Robot hoặc camera tự hành có thể theo dõi đàn vật nuôi, phát hiện con vật bị bệnh, kiểm tra nhiệt độ, hành vi ăn uống và hỗ trợ tự động cho ăn hoặc vệ sinh chuồng trại.

    \item \textbf{Sản xuất công nghiệp:} Robot tự hành trong nhà máy có thể vận chuyển linh kiện, cấp phát vật tư, kiểm tra dây chuyền, hỗ trợ lắp ráp và phối hợp với robot cánh tay để tăng năng suất.

    \item \textbf{Kho vận và phân phối:} Xe tự hành hoặc robot kho hàng có thể phân loại hàng hóa, đưa hàng đến đúng vị trí, kiểm kê tồn kho và hỗ trợ giao hàng chặng cuối.

    \item \textbf{Giám sát đô thị:} Drone tuần tra có thể quan sát giao thông, phát hiện ùn tắc, cháy nổ, vi phạm xây dựng hoặc các sự kiện bất thường ở khu vực đông người.

    \item \textbf{Cứu hộ và ứng phó khẩn cấp:} Trong cháy rừng, lũ lụt, động đất hoặc khu vực nguy hiểm, robot và drone có thể khảo sát hiện trường, tìm người mắc kẹt và gửi dữ liệu về trung tâm điều hành mà không đặt con người vào rủi ro trực tiếp.
\end{itemize}


\section*{Câu 13}
\addcontentsline{toc}{section}{Câu 13}
Triết lý của thiết bị tự hành là gì? Vì sao lại tồn tại, có cần thiết hay không? Vì sao tin vào tự hành?

Triết lý của thiết bị tự hành là chuyển một phần năng lực cảm nhận, suy luận và hành động từ con người sang máy móc, nhằm giúp hệ thống hoạt động an toàn, ổn định và hiệu quả hơn trong những môi trường phức tạp. Thiết bị tự hành không chỉ làm theo lệnh cố định, mà phải liên tục quan sát môi trường, hiểu tình huống, dự đoán rủi ro và chọn hành động phù hợp.

\begin{itemize}
    \item \textbf{Vì sao thiết bị tự hành tồn tại?} Chúng xuất hiện vì nhiều công việc trong thực tế có tính lặp lại, nguy hiểm, tốn sức hoặc cần phản ứng rất nhanh. Ví dụ, robot có thể vận chuyển hàng trong kho cả ngày, drone có thể bay vào vùng cháy rừng, còn xe tự hành có thể liên tục quan sát đường bằng nhiều cảm biến cùng lúc.

    \item \textbf{Có cần thiết hay không?} Thiết bị tự hành là cần thiết trong các hệ thống hiện đại vì chúng giúp tăng năng suất, giảm rủi ro cho con người và mở rộng khả năng giám sát ở những nơi con người khó tiếp cận. Tuy nhiên, chúng không thay thế hoàn toàn con người, mà nên đóng vai trò hỗ trợ, tự động hóa và giảm tải.

    \item \textbf{Vì sao có thể tin vào tự hành?} Có thể tin khi hệ thống được thiết kế với cảm biến đáng tin cậy, mô hình AI được kiểm thử kỹ, có cơ chế dự phòng, có khả năng phát hiện lỗi và có con người giám sát trong các tình huống quan trọng. Niềm tin vào tự hành không phải là tin tuyệt đối, mà là tin dựa trên kiểm chứng, dữ liệu và quy trình an toàn.

    \item \textbf{Giới hạn của tự hành:} Thiết bị tự hành vẫn có thể gặp lỗi khi môi trường quá phức tạp, dữ liệu cảm biến nhiễu, tình huống chưa từng gặp hoặc mô hình dự đoán sai. Vì vậy, cần có cơ chế dừng an toàn, ghi log, cảnh báo và cho phép con người can thiệp.
\end{itemize}

Tóm lại, thiết bị tự hành tồn tại vì xã hội cần các hệ thống có thể hỗ trợ con người làm việc nhanh hơn, an toàn hơn và bền bỉ hơn. Triết lý đúng không phải là ``máy thay con người'', mà là ``máy hỗ trợ con người'', trong đó con người vẫn giữ vai trò thiết kế mục tiêu, giám sát rủi ro và chịu trách nhiệm cuối cùng.